\section{Discussion and Limitations}
\label{sec:limitations}

\paragraph{One root per frame} Under DXVK and vkd3d-proton the presenting thread
is a translation thread, and a stall on the game thread can close in the
frame before the late one when frames are pipelined. The chain reaches the
game thread when the present thread's waits are ended by it, but the
prototype does not yet look back across a frame boundary.
\todo{Measure on Proton titles; consider identifying the frame thread from
the chain at present.}

\paragraph{One chain per group, one waker per wait} A group carries the chain
behind its longest interval only; its other waits may have had different
chains. A wait on several threads follows the one that finished last.

\paragraph{On-CPU kernel time is not ranked} Faults, reclaim and compaction are
carved out of on-CPU time as buckets rather than recorded as intervals, so
they form no group and never lead a record, however large
(\autoref{sec:eval-cases}, cases 8 and 9).

\paragraph{Only the last holder of a CPU} A hop names the task that last took
or held its CPU, not the one that held it longest. Once a background
program takes cores, the game's own threads start preempting each other,
and the last holder is sometimes one of them.

\paragraph{The gate does not yet use the resolution} The share gate weighs
blocked time including \textit{wait\_oncpu}. \autoref{sec:eval-gate}
shows the consequence and that subtracting it fixes it.

\paragraph{Not yet recorded} The requested duration of a sleep, so that an
oversleeping frame limiter is distinguishable from an intentional sleep; the
CPU each thread ran on, and whether other CPUs sat idle; how long an I/O
request queued in the block layer against how long the device took, and who
else had requests in flight; and a refinement of \bucket{block\_other} by
kernel stack (e.g.\ the ntsync driver~\cite{ntsync}, split-lock
mitigation~\cite{lwn-splitlock}). The case studies show where each is needed
(\autoref{sec:eval-cases}). \todo{Implement or scope out.}

\paragraph{Cost outside the game} \sys keeps state for every task on the
machine, because a chain can run through any process (a server, wineserver)
whose history is needed before it is known to matter. Every context switch
on the machine pays for it (\autoref{sec:eval-overhead}).
